\section{Method}
\label{sec:method}

We investigate how action-relevant structure is distributed across the components of a vision-language-action model by combining geometric analysis (RSA, CKA) with functional probes.
This section formalizes the model, data, representation extraction procedure, and the three complementary analysis tools.

%------------------------------------------------------------------
\subsection{Model and Data}
\label{sec:model_data}

\paragraph{Architecture.}
We study OpenVLA-7B~\cite{kim2024_openvla,karamcheti2024}, a generalist robot policy that chains three components: a SigLIP-SO400M vision encoder~\cite{zhai2023_siglip} (0.80B parameters, frozen during fine-tuning), an MLP projector that maps vision features into the language model's embedding space, and a Llama-2 7B language backbone (6.74B parameters) that autoregressively decodes 7-dimensional actions (6 DoF end-effector pose $+$ gripper opening).

\paragraph{Fine-tuning data.}
We use a checkpoint fine-tuned on the LIBERO-10 split~\cite{liu2023_libero}, which provides 10 tabletop manipulation tasks with 50 human demonstrations per task.

\paragraph{Evaluation benchmark.}
We evaluate on the LIBERO-Goal suite~\cite{liu2023_libero}, which shares initial states with LIBERO-10 but pairs them with different language-specified goals.
This separation ensures that the evaluation tasks test generalization to novel goals rather than memorization of training trajectories.

\paragraph{Untrained control.}
To isolate the contribution of LLM fine-tuning, we construct an untrained control by reinitializing the LLM parameters via \texttt{from\_config()}, yielding 6.74B random weights while preserving the frozen SigLIP encoder (0.80B pretrained weights) and the MLP projector architecture.
Any difference between trained and untrained conditions is therefore attributable to LLM training, not to the vision encoder.

%------------------------------------------------------------------
\subsection{Representation Extraction}
\label{sec:repr_extraction}

We extract representations from five sources that span the VLA processing pipeline and two external baselines:

\begin{enumerate}[label=(\alph*),leftmargin=*,nosep]
  \item \textbf{SigLIP+Projector}: the output of the MLP projector applied to SigLIP vision features, i.e., the representation immediately before LLM processing.
  \item \textbf{LLM image\_mean}: the mean activation across image token positions at each of the 32 LLM layers, capturing how distributed visual information evolves through the language backbone.
  \item \textbf{LLM last\_preaction}: the activation at the final token position before action decoding (the readout token) at each LLM layer, capturing how the model concentrates information for action prediction.
  \item \textbf{DINOv2 ViT-L/14}~\cite{oquab2024_dinov2}: the \texttt{cls\_token} from a vision-only self-supervised encoder, serving as a baseline without language grounding.
  \item \textbf{CLIP ViT-L/14}~\cite{radford2021_clip}: the mean of patch token activations (\texttt{patch\_mean}) from a vision-language contrastive model, serving as an external VL-pretrained baseline with no VLA training signal.
\end{enumerate}

\paragraph{Experimental conditions.}
We cross two factors in a $2 \times 2$ factorial design: \emph{training status} (trained vs.\ untrained LLM) and \emph{instruction presence} (with vs.\ without language instruction).
Sources (a)--(c) are extracted under all four conditions; sources (d) and (e) are condition-independent external baselines.
For the instruction factor, the with-instruction condition uses the LIBERO-Goal task description; the without-instruction condition provides an empty string.

%------------------------------------------------------------------
\subsection{Representational Similarity Analysis}
\label{sec:rsa}

We adopt Representational Similarity Analysis (RSA)~\cite{kriegeskorte2008} to quantify geometric alignment between neural representations and the action space.

\paragraph{RDM construction.}
For a set of $N$ observations, we compute the $N \times N$ representational dissimilarity matrix (RDM) for each representation source using pairwise Euclidean distances.
An analogous RDM is computed for the 7-dimensional action vectors.

\paragraph{RSA statistic.}
We measure the alignment between a representation RDM and the action RDM via Spearman rank correlation $\rho$ over the upper-triangular entries of each matrix.
Higher $\rho$ indicates that observations that are nearby in representation space also tend to require similar actions.

\paragraph{Task-level analysis.}
We compute RSA separately for each of the 10 LIBERO-Goal tasks, then aggregate across tasks.
This per-task procedure avoids conflating within-task structure with between-task variation.

\paragraph{Statistical inference.}
We construct 95\% confidence intervals via bootstrap resampling (10{,}000 resamples) over the observation pairs.
Pairwise comparisons between representation sources use bootstrap differences with the same resampling scheme.

%------------------------------------------------------------------
\subsection{Functional Action Probes}
\label{sec:probes}

RSA captures geometric alignment but does not assess whether action-relevant information is linearly decodable.
We complement RSA with functional probes following~\cite{alain2017}.

\paragraph{Probe architecture.}
We train two-layer MLP probes (Linear $\to$ ReLU $\to$ Linear; $d_\text{in} \to 256 \to 7$) to predict 7-D action vectors from frozen representations.
All representation weights are fixed; only the probe parameters are optimized.

\paragraph{Training protocol.}
Probes are trained with the Adam optimizer (learning rate $10^{-3}$, batch size 256) for up to 50 epochs with early stopping (patience $= 5$, monitoring validation loss).
Each task's data is split into training and validation sets.

\paragraph{Evaluation metric.}
We report mean $R^2$ across the 7 action dimensions, providing a normalized measure of variance explained that is comparable across representation sources with different dimensionalities.

\paragraph{Statistical comparison.}
We use bootstrap paired comparisons (10{,}000 resamples) to test whether $R^2$ differences between representation sources are statistically significant, reporting $p$-values and 95\% confidence intervals for all pairwise contrasts.

%------------------------------------------------------------------
\subsection{CKA Analysis}
\label{sec:cka}

As a complementary perspective, we employ Centered Kernel Alignment (CKA)~\cite{kornblith2019} alongside RSA.
While RSA captures local geometric alignment via rank-order preservation, CKA measures global representational similarity through kernel-based comparisons.

\paragraph{Two CKA variants.}
We compute both \emph{linear CKA}, which uses a linear kernel and is sensitive to second-order statistics, and \emph{RDM kernel CKA}, which operates on the same distance matrices as RSA but aggregates them differently.
Using both variants tests whether conclusions about representation structure are robust to the choice of similarity metric.

\paragraph{Dimension mismatch caveat.}
Linear CKA between high-dimensional representations (4{,}096-D for LLM layers) and the 7-D action space is subject to a known dimension mismatch: the linear kernel in the low-dimensional space may not capture the same structure as in the high-dimensional space.
RDM kernel CKA avoids this issue by operating on distance matrices of equal size regardless of the original feature dimensionality, providing a more reliable cross-dimensional comparison.
